% =====================================================================
% NeurIPS 2026 Paper Checklist
%
% Per the official NeurIPS 2026 instructions, the checklist must be
% included in the submitted PDF immediately after the technical
% appendix. It does NOT count toward the 9-page main-paper limit.
%
% Each item asks: did you do X? Answers must be one of [Yes] / [No] /
% [N/A], followed by a short justification. Do not leave items blank.
% =====================================================================

\section*{NeurIPS Paper Checklist}
\addcontentsline{toc}{section}{NeurIPS Paper Checklist}

\begin{enumerate}

\item {\bf Claims}\\
Do the main claims made in the abstract and introduction accurately
reflect the paper's contributions and scope?
\\\textbf{Answer: \answerYes{}}
\\\textbf{Justification:} The paper title and contributions frame
the work as a pre-registered audit and minimum reporting standard for
RAG faithfulness evaluation. The abstract and introduction explicitly
disclose: the matched-context CCS null; the small and
seed-sensitive scaled headline (no-refinement /
aggressive-refinement / gated-refinement, prior names HCPC-v$1$ /
HCPC-v$2$); the fixed-generation scorer fragility and
context-conditioned NLI sign-flip; the slice-dependent
scorer$-$human ranking between the typical-row calibration
($n{=}99$) and the targeted disagreement slice ($n{=}100$); the
threshold-transfer instabilities; and the cost-axis dominance of
the head-to-head. The contribution is the standard, not a new
method or leaderboard.

\item {\bf Limitations}\\
Does the paper discuss the limitations of the work performed by the
authors?
\\\textbf{Answer: \answerYes{}}
\\\textbf{Justification:} \S\ref{sec:limits} is grouped into three
categories: \emph{audited surface} (mostly short-answer SQuAD,
SQuAD-only matched-context, supplement-only $40$-question long-form
probe, mismatched Self-RAG harness); \emph{human evaluation}
(typical-row calibration $n{=}99$ with overlapping bootstrap CIs,
targeted disagreement slice $n{=}100$ adversarial-by-construction,
two raters plus adjudication on each); \emph{scope of replication}
(one local Qwen$2.5$ second-generator probe at $n{=}100$,
context-conditioned NLI rescoring on the $n{=}600$ balanced subset
not the full $7{,}500$). The section ends with the explicit
statement that these limits do not threaten the methodological
claim because the seven-axis decomposition is motivated by the
disclosure principle in \S\ref{sec:controlledrag}.

\item {\bf Theory Assumptions and Proofs}\\
For each theoretical result, does the paper provide the full set of
assumptions and a complete (and correct) proof?
\\\textbf{Answer: \answerNA{}}
\\\textbf{Justification:} The paper makes no formal theoretical
claims requiring proofs. All claims are empirical or definitional.

\item {\bf Experimental Result Reproducibility}\\
Does the paper fully disclose all the information needed to reproduce
the main experimental results of the paper to the extent that it
affects the main claims and/or conclusions of the paper (regardless of
whether the code and data are provided or not)?
\\\textbf{Answer: \answerYes{}}
\\\textbf{Justification:} \S\ref{sec:setup} specifies the pipeline,
embedder, indexer, reranker, generator, decoding settings, sample
sizes, seed list, and statistical protocol. The supplementary material
provides pre-specification, per-seed scaled-detail, and scorer-implementation
details. The artifact
ships per-query CSVs and scripts. \emph{Reproducibility scope}:
\texttt{run\_all\_analysis.sh} is end-to-end on a laptop in under
two minutes for the verification path (it regenerates every
analysis table in the main paper from the frozen per-query CSVs in
\texttt{results/revision/} and \texttt{human\_eval\_final/}).
Re-running the original \emph{generation} and \emph{scoring}
pipelines (matched-context generation, scaled audit, second-generator
probe, context-conditioned NLI rescoring) is a separate path that
needs local Ollama plus a CPU/MPS or free-tier GPU notebook and
takes hours; the supplementary compute notes document that path. The two-minute
verification path covers all numerical claims in the main paper.

\item {\bf Open access to data and code}\\
Does the paper provide open access to the data and code, with
sufficient instructions to faithfully reproduce the main experimental
results, as described in supplemental material?
\\\textbf{Answer: \answerYes{}}
\\\textbf{Justification:} The reviewer-facing anonymized artifact
(supplementary ZIP) contains all per-query CSVs, the human-evaluation
final files, analysis scripts, the pre-specification log, the source
trace, and a reproducibility README with verification commands.
Public release links and DOIs will be added at acceptance.

\item {\bf Experimental Setting/Details}\\
Does the paper specify all the training and test details (e.g., data
splits, hyperparameters, how they were chosen, type of optimizer,
etc.)?
\\\textbf{Answer: \answerYes{}}
\\\textbf{Justification:} The audit holds the entire RAG pipeline
fixed and varies only the audit axis under test. Embedder, indexer,
reranker, refinement policy, generator, and decoding settings are
listed in \S\ref{sec:setup}; sample sizes and seeds are listed for
every audited cell; the matched-context construction (similarity gap
$\leq 0.02$, CCS gap $\geq 0.20$) is given.

\item {\bf Experiment Statistical Significance}\\
Does the paper report error bars suitably and correctly defined or
other appropriate information about the statistical significance of
the experiments?
\\\textbf{Answer: \answerYes{}}
\\\textbf{Justification:} Every reported contrast is paired with a
$10{,}000$-resample bootstrap CI, paired Wilcoxon $p$-value where
applicable, Cohen's $d_z$ for paired contrasts, Wilson $95\%$ CIs for
binary rates, and bootstrap CIs for the cross-scorer correlations.
Multiple-comparison handling: only the matched-context test is
treated as the pre-specified primary; other paired contrasts are
descriptive across audit cells, as disclosed in \S\ref{sec:setup}.

\item {\bf Experiments Compute Resources}\\
For each experiment, does the paper provide sufficient information on
the computer resources (type of compute workers, memory, time of
execution) needed to reproduce the experiments?
\\\textbf{Answer: \answerYes{}}
\\\textbf{Justification:} The supplementary compute notes
specifies the local Ollama / Apple-silicon MPS / free-tier Kaggle T4
hardware, the per-fix runtime envelopes, and the analysis-only
verification path (\texttt{run\_all\_analysis.sh}, \textless 2 min on
any laptop). No paid APIs are required.

\item {\bf Code Of Ethics}\\
Does the research conform, in every respect, with the NeurIPS Code of
Ethics?
\\\textbf{Answer: \answerYes{}}
\\\textbf{Justification:} The work uses only public datasets, free or
local compute, no human-subjects data beyond two short-answer
support-judgment annotation passes by adult raters on public-dataset
content. No personally identifiable information, sensitive attributes,
or facial / biometric content is involved. No deceptive practices.
Annotator identities are not stored in the artifact.

\item {\bf Broader Impacts}\\
Does the paper discuss both potential positive societal impacts and
negative societal impacts of the work performed?
\\\textbf{Answer: \answerYes{}}
\\\textbf{Justification:} ControlledRAG's intended impact is to make
RAG faithfulness claims more honestly auditable, which reduces the
risk that deployed RAG systems are fielded based on under-identified
single-cell numbers (negative impact: false confidence in deployed
medical / legal RAG; the standard is meant to mitigate that). The
work is descriptive; it does not introduce a new generative model
that could itself be misused.

\item {\bf Safeguards}\\
Does the paper describe safeguards that have been put in place for
responsible release of data or models with a high risk for misuse?
\\\textbf{Answer: \answerNA{}}
\\\textbf{Justification:} The paper does not release a new generative
model, dataset of personally identifiable content, or other artifact
with high misuse risk. The released artifacts are evaluation tables,
scripts, and short-answer adjudicated annotation labels on public
data; these do not present such risks.

\item {\bf Licenses for existing assets}\\
Are the creators or original owners of assets (e.g., code, data,
models), used in the paper, properly credited and are the license and
terms of use explicitly mentioned and properly respected?
\\\textbf{Answer: \answerYes{}}
\\\textbf{Justification:} The paper cites SQuAD, PubMedQA, HotpotQA,
NaturalQuestions, TriviaQA, QASPER, MS-MARCO, and the open
embedder/cross-encoder/NLI models used. Each is used under its
published research-use license. No proprietary asset is used.

\item {\bf New Assets}\\
Are new assets introduced in the paper well documented and is the
documentation provided alongside the assets?
\\\textbf{Answer: \answerYes{}}
\\\textbf{Justification:} The new released assets are
(i)~the ControlledRAG seven-axis reporting standard
(\S\ref{sec:controlledrag}) plus a filled-in reporting-checklist
template
(\texttt{controlledrag\_reporting\_checklist.md},
distinct from this mandatory NeurIPS checklist),
(ii)~per-query CSVs for every audited cell,
(iii)~the typical-row calibration set ($n{=}99$),
(iv)~the targeted disagreement slice ($n{=}100$, two raters plus
adjudication, $74/26/0$ adjudicated distribution),
and (v)~the analysis scripts. Documentation:
\texttt{README\_REPRODUCE.md},
\texttt{SOURCE\_TRACE.md},
\texttt{artifact\_manifest.md},
\texttt{CLAIMS\_AUDIT.md},
and the supplementary source-trace section.

\item {\bf Crowdsourcing and Research with Human Subjects}\\
For crowdsourcing experiments and research with human subjects, does
the paper include the full text of instructions given to participants
and screenshots, if applicable, as well as details about compensation
(if any)?
\\\textbf{Answer: \answerYes{}}
\\\textbf{Justification:} The two adjudicated human-annotation arms
(n=99 typical-row calibration; n=100 targeted scorer-disagreement
evaluation) are documented with their exact label scheme, rater
instructions, decision codebook, and adjudication template; these
documents are shipped in the artifact under
\texttt{human\_eval\_final/}. No crowdsourcing platform was used; no
compensation was paid; the raters were not subject to risk.

\item {\bf Institutional Review Board (IRB) Approvals or Equivalent for Research with Human Subjects}\\
Does the paper describe potential risks incurred by study participants,
whether such risks were disclosed to the subjects, and whether
Institutional Review Board (IRB) approvals (or an equivalent approval/review
based on the requirements of your country or institution) were obtained?
\\\textbf{Answer: \answerNA{}}
\\\textbf{Justification:} The annotation tasks are short
support/faithfulness judgments on public-dataset content (no
personally identifiable information, no sensitive attributes, no
risk-of-harm content). Per common interpretations of the relevant
human-subjects regulations, this evaluation does not constitute
research with human subjects requiring IRB review. The annotators
were informed of the purpose of the labelling.

\item {\bf Declaration of LLM usage}\\
Does the paper describe the usage of LLMs if it is an important,
original, or non-standard component of the core methods in this
research? Note that if the LLM is used only for writing, editing, or
formatting purposes and does not impact the core methodology,
scientific rigorousness, or originality of the research, declaration
is not required.
\\\textbf{Answer: \answerYes{}}
\\\textbf{Justification:} A local Mistral-7B (via Ollama) is used as
both the audited generator and as the RAGAS-style local judge in the
metric-fragility audit. A local Qwen$2.5$-7B is used as the
second-generator replication. Their use, decoding settings, and the
fact that the RAGAS-style judge is itself an audit input (not a
ground-truth oracle) are disclosed in \S\ref{sec:setup} and
\S\ref{sec:metric_fragility}. LLMs are research artifacts in the
audited pipeline, not ground-truth oracles or co-authors.

\end{enumerate}
